\chapter*{Conclusion Générale et Perspectives}
\addcontentsline{toc}{chapter}{Conclusion Générale et Perspectives}
\label{chap:conclusion}

\noindent Ce projet de fin d'études a été réalisé chez la startup Uvey, en Tunisie. Il a
permis de concevoir, déployer et valider une solution répondant aux défis opérationnels
d'une équipe d'ingénierie gérant une architecture microservices en forte croissance. Au
delà d'une simple mise en œuvre de nouvelles technologies, ce projet a transfiguré les
pratiques de livraison logicielle et le mode opératoire en infrastructure de l'entreprise,
les rendant bien plus efficaces et reproductibles.

\medskip
\noindent Les objectifs du projet ont été totalement, ou pour le moins grandement,
atteints. Une infrastructure cloud-native sur Azure, répartie sur trois clusters \ac{AKS},
a été conçue et déployée via \textbf{Terraform}. Un pipeline \textbf{GitOps} complet a
été bâti, et une stack unifiée d'observabilité couvrant les logs, les métriques ainsi que
le tracing distribué a été mise en place. La migration de la base \textbf{PostgreSQL}
depuis \ac{GCP} vers Azure a été effectuée sans perte de données et sans interruption de
service. Un portail \textbf{Backstage} a été installé pour faciliter l'accès à l'ensemble
des outils de l'équipe.
\medskip

\noindent Ce projet a été un véritable apprentissage. La migration d'\textbf{Ingress NGINX}
vers \textbf{Application Gateway} avec \textbf{AGIC} a demandé de détailler les
annotations de chacun des microservices et de créer un environnement de test dédié. Lors
de la migration de \textbf{PostgreSQL}, des problèmes de compatibilité d'extensions entre
\textit{Cloud SQL} et \textit{Azure Flexible Server} ont nécessité des ajustements du
module \textbf{Terraform} avant de pouvoir reprendre le processus. Pour pallier aux
limitations de débit de l'API \textbf{ARM} d'Azure, nous sommes passés par un proxy
\textbf{Node.js} sur lequel nous avons implémenté un cache. D'un point de vue
organisationnel, la pratique de \textbf{Scrum} dans un contexte d'Infrastructure as Code
a révélé qu'il fallait prévoir des sprints avec plus de marge, ainsi qu'intégrer
explicitement la validation en environnement pré-prod dans la \textit{Definition of Done}.

\medskip
\noindent Pour dépasser ces challenges, nous sommes peu à peu convaincus qu'à ce jour,
aucune réussite sur du cloud public ne s'est affranchie de deux piliers fondamentaux~:
la \textbf{documentation architecturale} et l'\textbf{anticipation rigoureuse des
contraintes environnementales}. Ces deux dimensions sont tout aussi déterminantes que la
maîtrise technique pour réussir une transition vers le cloud.

%\bigskip

%\noindent \textbf{\large Perspectives :}

\medskip
\noindent Les travaux réalisés ouvrent plusieurs axes d'évolution naturels, structurés
autour de trois dimensions~:

\medskip
\noindent \textbf{Enrichissement de la plateforme développeur.} Le portail
\textbf{Backstage} pourrait être enrichi par l'intégration de \textit{scaffolders}
permettant aux équipes de provisionner elles-mêmes de nouveaux microservices conformes
aux standards architecturaux d'Uvey, réduisant ainsi davantage la friction entre
développement et infrastructure. Par ailleurs, l'adoption de \textbf{Karpenter} en
remplacement du \textit{Cluster Autoscaler} d'Azure permettrait un provisionnement des
n\oe{}uds \ac{AKS} plus réactif et économiquement optimisé, notamment grâce à la gestion
des instances Spot.

\medskip
\noindent \textbf{Renforcement de la sécurité et de l'observabilité.} Le déploiement
d'une stratégie réseau de type \textbf{zero-trust} avec \textbf{Cilium} ou \textbf{Azure
Network Policy} serait un prolongement naturel de l'existant avec \textit{Managed
Identity}, en allant jusqu'à contrôler chaque interconnexion entre les différents services
au sein des clusters. Du point de vue de l'observabilité, la généralisation du tracing
distribué vers un modèle \textbf{eBPF} permettrait de capter des signaux de performance
de manière non-intrusive, offrant une visibilité complète sur les échanges réseau
sous-jacents.

\medskip
\noindent \textbf{Conformité et certification.} À terme, la maturité de la plateforme
ouvrira la possibilité de viser une certification \textbf{ISO 27001} ou une attestation
\textbf{SOC 2}, des labels recherchés par les FinTechs pour médiatiser leur confiance
auprès de clients B2B et accélérer leur processus de vente sur des marchés régulés.